\documentclass[12pt, lettersize, journal]{article}
\usepackage[utf8]{inputenc}
\usepackage[indonesian]{babel}
\usepackage{amsmath}
\usepackage{amsfonts}
\usepackage{amssymb}
\usepackage{geometry}
\usepackage{listings}
\usepackage{xcolor}

\geometry{legalpaper, margin=1in}

% Konfigurasi Kode Program (Tab 2 Spasi)
\definecolor{codegreen}{rgb}{0,0.6,0}
\definecolor{codegray}{rgb}{0.5,0.5,0.5}
\definecolor{codepurple}{rgb}{0.58,0,0.82}
\definecolor{backcolour}{rgb}{0.95,0.95,0.92}

\lstdefinestyle{mystyle}{
    backgroundcolor=\color{backcolour},   
    commentstyle=\color{codegreen},
    keywordstyle=\color{magenta},
    numberstyle=\tiny\color{codegray},
    stringstyle=\color{codepurple},
    basicstyle=\ttfamily\footnotesize,
    breakatwhitespace=false,         
    breaklines=true,                 
    captionpos=b,                    
    keepspaces=true,                 
    numbers=left,                    
    numbersep=5pt,                  
    showspaces=false,                
    showstringspaces=false,
    showtabs=false,                  
    tabsize=2
}
\lstset{style=mystyle}

\title{\textbf{Metode Numerik: Penyelesaian SPL secara Iteratif (Gauss-Seidel)}}
\author{Catatan Kuliah Komputasi Numerik}
\date{\today}

\begin{document}

\maketitle

\section{Pendahuluan}
Berbeda dengan metode eliminasi langsung, metode **Iterasi Gauss-Seidel** menyelesaikan Sistem Persamaan Linear (SPL) secara aproksimasi berulang. Metode ini sangat efisien untuk SPL berukuran besar dengan matriks koefisien yang dominan secara diagonal (\textit{diagonally dominant}).

Kita memulai proses dengan menentukan tebakan awal untuk semua variabel, yang umumnya diinisialisasi dengan nilai nol:
\begin{equation}
    x_1^{(0)} = 0, \quad x_2^{(0)} = 0, \quad \dots, \quad x_n^{(0)} = 0
\end{equation}

\section{Algoritma dan Formulasi Matematika}
Untuk SPL dengan $n$ persamaan, kita mengisolasi setiap variabel $x_i$ pada persamaan ke-$i$. Rumus umum untuk iterasi ke-$(k+1)$ adalah:

\begin{equation}
    x_i^{(k+1)} = \frac{b_i - \sum_{j=1}^{i-1} a_{ij}x_j^{(k+1)} - \sum_{j=i+1}^{n} a_{ij}x_j^{(k)}}{a_{ii}}
\end{equation}

\subsection{Karakteristik Substitusi Balik Berulang}
Kunci utama Gauss-Seidel yang membedakannya dengan metode Jacobi adalah **penggunaan langsung nilai terbaru**. Begitu nilai $x_1^{(k+1)}$ didapatkan dari baris pertama, nilai baru tersebut langsung dipakai untuk menghitung $x_2^{(k+1)}$ di baris kedua pada iterasi yang sama. Proses ini terus berulang secara iteratif sampai selisih nilai antar-iterasi lebih kecil dari toleransi error yang ditentukan ($\epsilon$).

\section{Analisis Kasus Khusus ($m < n$)}
Jika jumlah persamaan lebih sedikit daripada jumlah variabel ($m < n$), metode iteratif seperti Gauss-Seidel akan menghadapi masalah serius:
\begin{itemize}
    \item \textbf{Matrix Underdetermined:} Kita tidak memiliki cukup persamaan untuk mengisolasi setiap variabel di sisi kiri rumus iterasi.
    \item \textbf{Divergensi:} Karena adanya variabel bebas yang tidak terikat oleh baris pivot, nilai iterasi kemungkinan besar akan terus membesar tanpa batas (\textit{diverge}) atau terjebak dalam lingkaran iterasi tak hingga tanpa pernah konvergen ke solusi manapun.
\end{itemize}

\section{Implementasi Code (Python)}
Berikut adalah implementasi metode Gauss-Seidel dengan tebakan awal nol dan ukuran tabulasi 2 spasi:

\begin{lstlisting}[language=Python, caption=Implementasi Iterasi Gauss-Seidel]
import numpy as np

def gauss_seidel(A, B, x_awal, tol=1e-5, max_iter=100):
  A = A.astype(float)
  B = B.astype(float)
  x = x_awal.copy().astype(float)
  n = len(B)
  
  for k in range(max_iter):
    x_lama = x.copy()
    
    for i in range(n):
      # Hitung sigma untuk variabel yang sudah diupdate dan belum
      sigma = 0.0
      for j in range(n):
        if j != i:
          sigma += A[i][j] * x[j]
          
      # Ambil nilai terbaru langsung (Back Substitution internal)
      x[i] = (B[i] - sigma) / A[i][i]
      
    # Cek konvergensi menggunakan norm
    if np.linalg.norm(x - x_lama, ord=np.inf) < tol:
      print(f"Konvergen pada iterasi ke-{k+1}")
      return x
      
  print("Metode tidak konvergen dalam batas maksimum iterasi")
  return x

# Contoh Kasus
A = np.array([[4, 1, -1], [2, 7, 1], [1, -3, 12]])
B = np.array([3, 19, 31])
# Tebakan awal semua variabel = 0
tebakan_awal = np.array([0, 0, 0])

solusi = gauss_seidel(A, B, tebakan_awal)
print("Solusi:", solusi)
\end{lstlisting}

\end{document}